\documentclass[12pt]{article}
\usepackage[utf8]{vietnam}
\usepackage[margin=1in]{geometry}
\usepackage{titlesec}
\usepackage{array}
\usepackage{hyperref}
\usepackage{xcolor}

\title{\textbf{BÁO CÁO BUỔI 2: KIẾN TRÚC REST VÀ HTTP FUNDAMENTALS}}
\author{Kiến trúc hướng dịch vụ (SOA)}
\date{2026}

\begin{document}

\maketitle

\section{Tổng quan về REST}
Hệ thống World Wide Web giống như một thành phố, trong đó \textbf{HTTP} là hệ thống đường xá, còn \textbf{REST} là bộ luật giao thông giúp mọi thứ vận hành trơn tru.

\section{6 Nguyên tắc vàng của kiến trúc REST}
Để một hệ thống được công nhận là RESTful, nó phải tuân thủ các ràng buộc sau:

\begin{itemize}
    \item \textbf{Client-Server:} Tách biệt giữa giao diện và dữ liệu, cho phép phát triển độc lập.
    \item \textbf{Stateless (Phi trạng thái):} Mỗi request phải chứa đầy đủ thông tin để Server hiểu được mà không cần lưu lại ngữ cảnh từ các request trước.
    \item \textbf{Cacheable (Có thể lưu bộ nhớ đệm):} Phản hồi phải tự định nghĩa khả năng lưu cache để tăng hiệu năng.
    \item \textbf{Uniform Interface (Giao diện thống nhất):} Mọi tài nguyên được định danh qua URI và tương tác qua các phương thức chuẩn.
    \item \textbf{Layered System (Hệ thống phân lớp):} Client không cần biết mình đang kết nối trực tiếp với Server hay qua các lớp trung gian (Proxy, Load Balancer).
    \item \textbf{Code on Demand (Tùy chọn):} Server có thể gửi mã thực thi (như JavaScript) về phía Client.
\end{itemize}

\section{HTTP Fundamentals}

\subsection{Các phương thức HTTP (Methods)}
\begin{center}
\begin{tabular}{|m{2cm}|m{5cm}|m{6cm}|}
    \hline
    \textbf{Method} & \textbf{Mục đích} & \textbf{Tính chất (Idempotent)} \\
    \hline
    GET & Lấy dữ liệu & Có (Gọi nhiều lần kết quả không đổi) \\
    \hline
    POST & Tạo mới tài nguyên & Không (Tạo bản ghi mới mỗi lần gọi) \\
    \hline
    PUT & Ghi đè toàn bộ & Có (Thay thế hoàn toàn tài nguyên) \\
    \hline
    PATCH & Cập nhật một phần & Thường không (Chỉ sửa trường cụ thể) \\
    \hline
    DELETE & Xóa tài nguyên & Có (Xóa rồi thì vẫn là đã xóa) \\
    \hline
\end{tabular}
\end{center}

\subsection{Mã trạng thái (Status Codes)}
\begin{itemize}
    \item \textbf{\color{blue}2xx (Success):} Thành công (Ví dụ: 200 OK, 201 Created).
    \item \textbf{\color{orange}3xx (Redirection):} Chuyển hướng (Ví dụ: 301 Moved Permanently).
    \item \textbf{\color{red}4xx (Client Error):} Lỗi phía người dùng (Ví dụ: 400 Bad Request, 404 Not Found).
    \item \textbf{\color{purple}5xx (Server Error):} Lỗi phía hệ thống (Ví dụ: 500 Internal Server Error).
\end{itemize}

\section{Mô hình trưởng thành Richardson (Maturity Model)}
Đánh giá mức độ RESTful của một API qua 4 cấp độ:
\begin{description}
    \item[Level 0:] Dùng HTTP như đường dẫn truyền tin (thường chỉ dùng POST).
    \item[Level 1:] Sử dụng Tài nguyên (Resources) với các URL riêng biệt.
    \item[Level 2:] Sử dụng đúng các Động từ HTTP và Mã trạng thái.
    \item[Level 3:] \textbf{HATEOAS} (Hypermedia as the Engine of Application State).
\end{description}

\end{document}
